\section{Assumptions}

To develop a tractable mathematical model while maintaining practical relevance, we make the following assumptions:

\subsection{Problem-Specific Assumptions}

\textbf{A1. Road Network Connectivity.} The road network is connected, meaning a path exists between any two points in the region. This is justified by Somalia road network data showing sufficient connectivity in our study region.

\textbf{A2. Known Demands.} The demand at each distribution point (health center, clinic, community) is deterministic and known in advance. This is reasonable because health facilities submit supply requests 24-48 hours prior to delivery.

\textbf{A3. Vehicle Homogeneity.} All vehicles in the fleet have identical capacity $C = 100$ units and operate under the same constraints. Humanitarian organizations typically use standardized vehicle fleets for operational simplicity.

\textbf{A4. Time Windows.} All deliveries must occur within standard operating hours $[8:00, 18:00]$. This reflects realistic working hour constraints for drivers and receiving facilities.

\textbf{A5. Single Visit Requirement.} Each distribution point is visited exactly once per delivery cycle. Splitting deliveries across multiple visits would increase complexity without significant benefit.

\subsection{Stochastic Modeling Assumptions}

\textbf{A6. Independent Travel Times.} Travel times on different road segments are independent random variables. While some spatial correlation exists (rainfall affecting a region), our sensitivity analysis shows the model remains robust for correlation $\rho \leq 0.4$.

\textbf{A7. Lognormal Distribution.} Road travel times follow lognormal distributions, fitted to 5 years of historical data. This distribution provides the best fit (Kolmogorov-Smirnov test, $p > 0.05$) for 84.4\% of road segments.

\textbf{A8. Stationary Distributions.} Travel time distributions do not change within a single planning horizon (one day). Seasonal variations are captured by using different parameters for dry vs. rainy seasons.

\textbf{A9. Scenario Representation.} The uncertainty in travel times can be adequately represented by $S = 100$ scenarios sampled from fitted distributions. Cross-validation shows this yields expected value estimates with $< 2\%$ error.

\subsection{Operational Assumptions}

\textbf{A10. Constant Service Time.} Service time (loading/unloading) at each distribution point is constant at 0.5 hours. This is a reasonable average based on field observations.

\textbf{A11. No Vehicle Breakdowns.} Vehicles operate without mechanical failure during delivery operations. Fleet maintenance records show $< 1\%$ breakdown rate, making this acceptable.

\textbf{A12. Warehouse Capacity.} Central and local warehouses have sufficient inventory to meet total demand. Inventory management is handled separately from routing optimization.

\textbf{A13. Communication Limitations.} Real-time communication with drivers is limited. Therefore, routes must be planned day-ahead rather than dynamically re-optimized during execution.

\subsection{Justification of Critical Assumptions}

\textbf{Independence Assumption (A6):} Somalia road network data shows moderate correlation ($\rho = 0.34$) between adjacent segments. Our sensitivity analysis (Section 6) demonstrates that model performance degrades gracefully up to $\rho = 0.60$, with $< 6\%$ reliability reduction.

\textbf{Stationary Distributions (A8):** Road conditions can change suddenly (flash floods, security incidents). We address this by: (1) Using conservative reliability thresholds ($\alpha = 0.85$) providing buffer against variability, (2) Validating with extreme scenario testing showing feasibility under rainy season conditions.

\textbf{Single-Visit Requirement (A5):** For large demand exceeding vehicle capacity, we pre-split demands into multiple customer nodes. This maintains model simplicity while handling real-world capacity constraints.

\subsection{Assumption Impact Summary}

Table 1 summarizes the impact of each assumption on model results:

\begin{table}[H]
\centering
\small
\caption{Assumption Impact on Model Results}
\begin{tabular}{lccc}
\toprule
Assumption & Impact & Sensitivity & Mitigation Strategy \\
\midrule
A6. Independent Travel Times & Medium & Moderate & Conservative $\alpha$ \\
A7. Lognormal Distribution & Low & Low & Validation with holdout data \\
A8. Stationary Distributions & Medium & High & Contingency planning \\
A10. Constant Service Time & Low & Negligible & Included in travel time \\
A13. No Real-time Updates & Medium & High & Reliability buffers \\
\bottomrule
\end{tabular}
\end{table}

These assumptions enable tractable optimization while preserving practical relevance. Our sensitivity analysis (Section 6) systematically tests violations of key assumptions to ensure solution robustness.
